\section{Conclusion and Future Research}\label{h1:discuss}

In this work, we investigated the problem of learning to play zero-sum games against an arbitrary adversary under bandit feedback. We introduced the PSMR metric to quantify the learner's performance against the safety value a deterministic learner can achieve against the worst-case adversary. We provide comprehensive results, tackling both uninformed and informed feedback models and covering both normal-form and bilinear games.

Our results point toward several promising future directions. First,
as discussed, the regret bounds of Tsallis-INF and \PureUCB{} are incomparable. 
    Is it possible to achieve both simultaneously in the informed setting?
    Taking one step further, what are the instance-optimal bounds for this setting? 
    
    
Second, while our results for general bilinear games are analogous to those for normal-form games, it has been shown that in the special case of stochastic linear bandits ($|\calY|=1$), such bounds are not instance-optimal~\citep{lattimore2017end}.
    Therefore, understanding what fundamental quantity characterizes the difficulty of a general bilinear game is a promising direction.
    

Finally, our results are restricted to $\psmr$. While it coincides with $\NR$ when a PSNE exists, investigating how to achieve similar results for $\NR$ in the general case is an interesting direction.

